\documentclass[a4paper,12pt]{article}

% === Кодировки и шрифты для русского языка ===
\usepackage[T2A]{fontenc}
\usepackage[utf8]{inputenc}
\usepackage[russian]{babel}
\usepackage{indentfirst}

% === Математика и графика ===
\usepackage{amsmath, amssymb, amsfonts, mathtools, bm}
\usepackage{graphicx, float, booktabs, multirow}
\usepackage{adjustbox} % Для широких таблиц
\usepackage{array}     % Для контроля ширины колонок

% === Форматирование ===
\usepackage[left=2cm, right=2cm, top=2cm, bottom=2cm]{geometry} 
\usepackage{setspace}
\onehalfspacing

% === Ссылки и библиография ===
\usepackage{cite, url, hyperref}
\hypersetup{colorlinks=true, linkcolor=black, citecolor=blue, urlcolor=blue}

% === Нумерация формул по разделам ===
\numberwithin{equation}{section}

\newcommand{\E}{\mathbb{E}}
\newcommand{\BER}{\mathrm{BER}}
\newcommand{\SNR}{\mathrm{SNR}}

\title{\textbf{Частотное разнесение в OFDM-системах через управляемое наложение спектральных реплик: Архитектура ПСИР для использования зон Найквиста}}
\author{Рычков~Е.~Н.$^{1}$ \\ $^{1}$Сибирский Федеральный университет, \texttt{meurchmusic@gmail.com}}
\date{}

\begin{document}
\maketitle

\begin{abstract}
Предложена архитектура трансивера сигналов OFDM, реализующая принцип использования алиасинга и зон Найквиста в концепции
спектрально-фазового разнесения (Spectral-Phase Diversity) на основе детерминированного использования спектральных реплик, возникающих при дискретизации сигнала цифро-аналоговыми преобразователями (ЦАП). Традиционные методы рассматривают алиасинг как необратимую деградацию сигнала; предлагаемый подход трансформирует эти компоненты в независимые каналы связи, что можно рассматривать в контексте новой степени свободы для оборудования с ограниченной полосой частот. Метод, обозначаемый далее как ПСИР (Phase-Shifted Interleaved Receiving), использует сдвиг тактовой фазы $\Delta t = T_s/M$ между интерливинговыми АЦП для создания предсказуемой ортогональной структуры смеси, что позволяет обратимо выделить реплики с помощью унитарного преобразования (матрицы Адамара для $M=2$) и применить максимально правдоподобное комбинирование (MRC). Ключевое преимущество метода заключается в реализации частотного разнесения без расширения полосы передаваемого сигнала и без дополнительных ВЧ-трактов за счёт утилизации спектральных реплик ЦАП (при условии обеспечения входной аналоговой полосы приёмника $[f_s, f_s+B]$ и занятия дополнительного спектра в зонах Найквиста). ПСИР комплементарен технологиям MIMO: поскольку методы используют ортогональные ресурсы (пространство vs. зоны Найквиста), число независимых ветвей разнесения перемножается ($L_{\text{total}} = L_{\text{freq}} \cdot L_{\text{space}}$), позволяя получать надёжность многоантенных систем на компактных устройствах. Проведена численная верификация методом Монте-Карло для базового случая $M=2$. Показано, что схема удваивает конструктивный порядок разнесения с $L=1$ (локальный наклон $d = 0.968$) до $L=2$ ($d = 1.872$ в окне 14--24~дБ; $d = 2.102$ в окне 22--30~дБ), обеспечивая выигрыш 5.24~дБ (5.26~дБ аналитически с учётом шума LS-оценки) по отношению сигнал/шум при BER $10^{-2}$ относительно аппаратного контроля. Продемонстрирована ортогональность метода пространственному разнесению: гибрид ПСИР + SIMO $1{\times}2$ достигает асимптоты $L=4$, а гибрид ПСИР + MIMO $2{\times}2$ Alamouti достигает асимптоты $L=8$ с выигрышем 2.43~дБ (2.45~дБ аналитически) поверх чистого Alamouti при BER $10^{-4}$. Все ключевые утверждения подтверждены машинными тестами, встроенными в код симуляции. Обсуждается возможность адаптации архитектуры к условиям ограниченного спектра путем уплотнения зон Найквиста. Метод представляет интерес для энергоэффективных систем (IoT, спутниковая связь, спецсвязь), где приоритет отдается надежности и миниатюризации оборудования над максимальной спектральной эффективностью.

\vspace{0.3cm}
\noindent\textbf{Ключевые слова}: частотное разнесение, спектральные реплики, зоны Найквиста, интерливинговые АЦП (TI-ADC), управляемый алиасинг, матрица Адамара, когерентное комбинирование (MRC), Direct RF Sampling, ПСИР, гибридные системы MIMO.
\end{abstract}

\section{Введение}

Обеспечение высокой надежности передачи данных в условиях многолучевого распространения является одной из ключевых задач современной радиотехники. Традиционные методы борьбы с замираниями включают использование временных, частотных и пространственных ресурсов. 

Пространственное разнесение (MIMO) эффективно повышает пропускную способность и надежность, однако требует множества отдельных высокочастотных трактов и антенн, что значительно увеличивает стоимость, габариты и энергопотребление устройства \cite{proakis,goldsmith}. Частотное разнесение, реализуемое через дублирование данных в широкой полосе пропускания или на разнесённых поднесущих \cite{abdelgader}, часто ограничено дефицитом радиочастотного спектра: расширение полосы для получения независимых каналов ведет к снижению спектральной эффективности и усложнению регуляторного соответствия \cite{goldsmith}.

\subsection{Проблематика и идея работы}
В цифровой радиотехнике эпохи Direct RF Sampling (прямой оцифровки ВЧ-сигналов) процесс дискретизации неизбежно порождает спектральные копии (реплики) сигнала в зонах Найквиста. Инженерная практика традиционно направлена на подавление этих копий с помощью аналоговых фильтров нижних частот (LPF) на выходе ЦАП, рассматривая их как источник искажений и потерь энергии.

Идея данной работы состоит в использовании спектральных копий не как помехи, подлежащей фильтрации, а как дополнительного измерения пространства состояний канала. Мы показываем, что при условии точной синхронизации тактовых сигналов приемника эти копии становятся статистически независимыми ветвями канала связи (при выполнении условия $f_s > B_c$). Это позволяет достичь эффектов, аналогичных MIMO, но на одном компактном устройстве с минимальным набором активных элементов.

Данный подход предполагает компромисс: он снижает спектральную эффективность (бит/с/Гц) из-за дублирования информации в разных зонах Найквиста, но существенно упрощает аппаратную часть передатчика и приемника, устраняя необходимость в множественных гетеродинах и антеннах. Для приложений с жесткими ограничениями по массе, габаритам и энергопотреблению (IoT, микро-спутники, носимая электроника) такой обмен является стратегически выгодным. Более того, метод дополняет технологию MIMO, позволяя строить гибридные системы с кратным увеличением надежности.

\subsection{Факторы, препятствовавшие развитию подхода}
Несмотря на известность явления наложения спектров, систематическое использование спектральных образов для гарантированного частотного разнесения не получило широкого распространения из-за трех причин:

\begin{enumerate}
    \item \textbf{Терминологическая фрагментация.} Концепции \textit{bandpass sampling} (полосовая дискретизация) \cite{vaughan,harris,rajasekaran}, \textit{Nyquist folding} (сворачивание спектра) \cite{fudge,murray} и \textit{Direct RF sampling} исследуются разными сообществами отдельно друг от друга. Специалисты по проектированию АЦП фокусируются на подавлении паразитных копий, а специалисты по цифровой обработке сигналов рассматривают алиасинг как необратимую потерю информации. Ни одна из групп не ставила целью обратимое использование наложения спектров для разнесения.
    
    \item \textbf{Недоопределенность системы при одном АЦП.} В классических схемах фаза наложения реплик детерминирована, но неизвестна приемнику из-за неопределенности начальной фазы несущей и канала, что делает систему недоопределенной (одно уравнение, две неизвестные). Использование технологий точного фазирования тактовых генераторов (Time-Interleaved ADC, TI-ADC) позволило создать предсказуемую фазовую структуру смеси \cite{walden,tsai,rajeswaran}, пригодную для математического обращения. Вопросы калибровки таких систем остаются актуальными \cite{elbornsson}.
    
    \item \textbf{Приоритет спектральной эффективности.} В коммерческих сетях (LTE, 5G NR) основным показателем является бит/с/Гц. Предлагаемая архитектура занимает нишу приложений, где спектр доступен или избыточен (ISM, SatCom, спецсвязь), а жесткие ограничения накладываются на массу, габариты и потребление электроэнергии приемника.
\end{enumerate}

\subsection{Основные вклады работы}
\begin{enumerate}
    \item Предложена и формализована архитектура приёмника (ПСИР), реализующая принцип спектрально-фазового разнесения, использующая детерминированное разделение спектральных реплик ЦАП для создания виртуального частотного разнесения без увеличения числа антенн и высокочастотных трактов.
    \item Разработан алгоритм изоляции независимых каналов связи с помощью унитарной матрицы Адамара после фазовой коррекции на уровне каждой поднесущей. Доказана необходимость соблюдения строгого порядка операций: разделение должно предшествовать оценке канала.
    \item Выполнена строгая численная верификация архитектуры методом Монте-Карло с идентификацией порядка разнесения по форме всей кривой ($\hat{L}=2$, локальный наклон $d = 1.872$ в окне 14--24~дБ), подтверждена эквивалентность классическому частотному разнесению при идентичном аппаратном бюджете, а также устойчивость к типовым неидеальностям TI-ADC методом общих случайных чисел (макс. $|\Delta\BER/\BER| = 0.25\%$ в диапазоне 0--30~дБ и $+0.23\%$ на тесте отсутствия порога ошибок при 30~дБ) \cite{ieee_std,adi_mt008}. В модель встроены машинные тесты, проверяющие вещественный тракт, нормировку мощности, порядок оценки канала и энергетический бюджет.     \item Показана совместимость и мультипликативная ортогональность метода с технологиями пространственного разнесения при равном числе АЦП: гибриды ПСИР + SIMO $1{\times}2$ ($\hat{L}=4$) и ПСИР + MIMO $2{\times}2$ Alamouti ($\hat{L}=8$) достигают теоретических пределов \cite{proakis} с выигрышем 2.43~дБ (2.45~дБ аналитически с учётом потерь LS-оценки~\eqref{eq:gamma_eff}) поверх чистого Alamouti при BER $10^{-4}$ (табл.~\ref{tab:results}). Обсуждается теоретическая масштабируемость метода на $M$ зон Найквиста и способы оптимизации занятой полосы.
\end{enumerate}

\section{Теоретические основы и модель системы}

\subsection{Формирование спектральных реплик и проблема наложения спектров}

Пусть передаваемый сигнал представляет собой OFDM-символ длительностью $T_u$. Комплексная огибающая этого символа $x(t)$ формируется как сумма модулированных поднесущих:
\begin{equation}
x(t) = \sum_{k=0}^{K-1} X[k] e^{j 2 \pi (k \Delta f) t}, \quad t \in [0, T_u],
\label{eq:ofdm_baseband}
\end{equation}
где $X[k]$ — комплексные амплитуды поднесущих, $\Delta f = 1/T_u$ — шаг сетки частот, $K$ — число активных поднесущих. Полоса сигнала равна $B = K \Delta f$.

При формировании реального проходного сигнала цифро-аналоговым преобразователем (ЦАП) с частотой обновления $f_s = 1/T_s$ возникает периодическое продолжение спектра в виде спектральных реплик (образов), центрированных на частотах $m f_s$, где $m \in \mathbb{Z}$. Каждая такая реплика несет в себе ту же самую информационную структуру $x(t)$, но смещенную в частотной области.

Рассмотрим две основные реплики, попадающие в разные зоны Найквиста приемника: первую ($N=1$, центр $f_0$) и вторую ($N=2$, центр $f_0 + f_s$). Эти реплики проходят через общий радиоканал с импульсной характеристикой $g(\tau)$. Поскольку частотный разнос между центрами реплик равен $f_s$, а ширина полосы каждой реплики равна $B$, условие статистической независимости замираний выполняется, если $f_s > B_c$, где $B_c$ — когерентная ширина полосы канала \cite{goldsmith}. В этом случае частотные отклики канала в окрестностях $f_0$ и $f_0+f_s$ можно рассматривать как независимые случайные величины.

В рамках символьной модели (после операции БПФ на приеме) воздействие канала на каждую поднесущую $k$ первой и второй реплик описывается комплексными коэффициентами передачи $\alpha_1[k]$ и $\alpha_2[k]$ соответственно. Тогда наблюдаемые сигналы в частотной области могут быть представлены как суперпозиция этих компонентов. Переход к вещественному времени для анализа интерференции дает следующую запись полного сигнала на входе приемника (до дискретизации АЦП, но после прохождения канала):
\begin{equation}
s(t) = \Re \left\{ \tilde{s}_1(t) e^{j 2 \pi f_0 t} + \tilde{s}_2(t) e^{j 2 \pi (f_0 + f_s) t} \right\},
\label{eq:real_signal}
\end{equation}
где $\tilde{s}_1(t)$ и $\tilde{s}_2(t)$ — комплексные огибающие первой и второй реплик после прохождения канала. В идеализированном случае плоского канала внутри полосы OFDM и независимости зон Найквиста, $\tilde{s}_i(t) \approx \alpha_i x(t)$. Таким образом, для каждого элемента вектора поднесущих $X[k]$ эффективные каналы $\alpha_1[k]$ и $\alpha_2[k]$ считаются независимыми рэлеевскими величинами.

\subsubsection*{Математическая природа наложения спектров}
Главная трудность, препятствующая использованию реплик для разнесения в классических схемах, заключается в природе их наложения при дискретизации одним АЦП. Согласно теории полосовой дискретизации \cite{vaughan,harris,vaidyanathan}, спектральный образ, расположенный вокруг частоты $f_0+f_s$, при попадании в базовую зону Найквиста приобретает фазовый множитель, зависящий от абсолютной частоты компонента:
\begin{equation}
S_{alias}[f] \propto S[f - f_s] \cdot e^{-j 2 \pi (f-f_s) T_s} = S[f - f_s] \cdot e^{-j 2 \pi f T_s},
\label{eq:alias_phase}
\end{equation}
где $T_s = 1/f_s$. Поскольку фаза $e^{-j 2 \pi f T_s}$ меняется линейно с частотой $f$, различные спектральные компоненты первой и второй реплик складываются с произвольными относительными фазами. Результатом является нелинейная интерференция, которую невозможно устранить простой фильтрацией без значительной потери энергии полезного сигнала. Именно поэтому традиционный подход считает алиасинг необратимой деградацией.

Предлагаемый нами метод обходит это ограничение, используя два АЦП со сдвигом тактов $\Delta t = T_s/2$. Этот сдвиг трансформирует случайную фазовую зависимость \eqref{eq:alias_phase} в детерминированную разность фаз $\pi$ между репликами, что позволяет обратить процесс смешивания с помощью ортогонального преобразования (матрицы Адамара).

\textbf{Важное условие эффективности:} Необходимо выполнение неравенства $f_s > B_c$. Это гарантирует статистическую независимость $\alpha_1$ и $\alpha_2$ \cite{goldsmith}. Если это условие нарушено, реплики становятся коррелированными, и эффект разнесения исчезает.

\subsection{Почему обычный алиасинг не дает разнесения? Анализ интерференции}
Если бы мы использовали один АЦП, обе реплики наложились бы друг на друга в первую зону Найквиста. Пусть $Y_{single}[k]$ — наблюдение на выходе единственного АЦП. Согласно свойствам аддитивности шума и линейности дискретизации, сигнал принимает вид:
\begin{equation}
Y_{single}[k] = (\alpha_1[k] + \beta_k \alpha_2[k]) x[k] + n_{combined}[k],
\label{eq:single_adc}
\end{equation}
где $\beta_k = e^{j \theta_k}$ — случайный фазовый множитель, определяемый соотношением несущей частоты и частоты дискретизации ($f_0/f_s$), а $n_{combined}[k]$ — суммарный шум. Поскольку $\alpha_1$ и $\alpha_2$ являются независимыми комплексными гауссовскими случайными величинами (CGRV), их сумма с случайной фазой также является CGRV с удвоенной дисперсией мощности. Однако, с точки зрения теории разнесения, это эквивалентно наличию \textit{одной} ветви с повышенной средней мощностью, но с тем же порядком распределения замираний (Rayleigh). Порядок разнесения $d$ определяется числом \textit{статистически независимых} ветвей, доступных для когерентного объединения. В уравнении \eqref{eq:single_adc} информация о $\alpha_1$ и $\alpha_2$ неразделимо смешана; невозможно оценить каждый канал отдельно для применения MRC. Следовательно, система ведет себя как одноветвлевая ($d=1$), несмотря на то, что физически сигнал прошел через два разных участка спектра. Отсутствие возможности декоррелировать каналы делает невозможным получение diversity gain выше единицы \cite{proakis}.

Более того, в области низких значений отношения сигнал/шум такая конструкция страдает от эффекта глубокого замирания: если векторы каналов $\alpha_1$ и $\alpha_2$ противоположны по фазе ($\arg(\alpha_1) \approx \arg(\alpha_2) + \pi$), они взаимно уничтожаются, приводя к резкому росту вероятности ошибки. Алиасинг здесь выступает как помеха, разрушающая информацию и делающую невозможным простое усреднение ветвей.

\subsection{Прием с интерливинговыми АЦП: детерминированная фаза как ключ к обратимости}
На приеме используются два идентичных АЦП, работающих с общей частотой дискретизации $f_s$, но со сдвигом момента выборки $\Delta t = T_s / 2$, где $T_s = 1/f_s$. Такой режим соответствует технологии Time-Interleaved ADC (TI-ADC) \cite{walden,rajeswaran,tsai}. Структура TI-ADC представляет собой частный случай многофазных обрабатывающих цепочек, подробно изученных в теории многоскоростных систем \cite{vaidyanathan}.

Для поднесущей с нормированной частотой $\nu_k = f_k / f_s$ сигналы на выходах АЦП в частотной области (после БПФ) принимают вид. Рассмотрим влияние сдвига тактов на спектральную компоненту $X[k]$. Задержка во временной области $t \to t - \tau$ соответствует умножению на $e^{-j 2 \pi f \tau}$ в частотной области \cite{oppenheim}. Для второго АЦП эффективная задержка относительно первого составляет $T_s/2$. Однако, поскольку вторая реплика находится в смещенной зоне Найквиста (сдвиг частоты на $f_s$), ее фаза меняется знакопеременно.

С учётом нормировки полной передаваемой мощности между репликами ($\rho_1 + \rho_2 = 1$, в базовом режиме $\rho_1 = \rho_2 = 0.5$) наблюдения в частотной области записываются как:
\begin{align}
Y_1[k] &= \left( \sqrt{\rho_1}\alpha_1[k] + \sqrt{\rho_2}\alpha_2[k] \right) x[k] + n_1[k], \label{eq:y1} \\
Y_2[k] &= e^{j \pi \nu_k} \left( \sqrt{\rho_1}\alpha_1[k] - \sqrt{\rho_2}\alpha_2[k] \right) x[k] + n_2[k], \label{eq:y2}
\end{align}

где $n_1[k]$ и $n_2[k]$ — аддитивный белый гауссов шум с дисперсией $\sigma_n^2$. Оценка спектра и статистические свойства шума подробно рассмотрены в классических работах по спектральному анализу \cite{black}.

\textbf{Вывод множителя $e^{j \pi \nu_k}$:}
Фаза, накопленная за время $T_s/2$ для частоты $f_k$, равна $\phi = 2 \pi f_k (T_s/2) = \pi (f_k/f_s) = \pi \nu_k$. Знак минус перед $\alpha_2$ обусловлен фазовым набегом $e^{j \pi f_s / f_s} = -1$ для второй реплики при сдвиге выборки на $T_s/2$, что позволяет алгебраически разделить компоненты без остаточной интерференции \cite{vaughan}. Именно эта \textit{детерминированная} зависимость позволяет восстановить исходные компоненты $\alpha_1 x$ и $\alpha_2 x$ без остаточной интерференции.

\textbf{Линейность и отсутствие аппаратного гашения.} Важно подчеркнуть, что система уравнений \eqref{eq:y1}-\eqref{eq:y2} является линейной. Даже если в конкретный момент времени происходит деструктивная интерференция в первом канале ($\alpha_1 \approx -\alpha_2 \implies Y_1 \approx 0$), во втором канале наблюдается конструктивная интерференция ($\alpha_1 - \alpha_2 \approx 2\alpha_1 \implies |Y_2|$ максимален). Информация не теряется в аналоговом тракте, а перераспределяется между каналами наблюдения. Обратное преобразование (матрица Адамара) позволяет алгебраически выделить оригинальные компоненты $\alpha_1 x$ и $\alpha_2 x$ с точностью до шума, предотвращая эффект «нулевого приема», характерный для одноканальных схем с алиасингом.

\subsection{Алгоритм разделения и комбинирования}
Процесс демодуляции включает следующие шаги, порядок которых критически важен:

\begin{enumerate}
    \item \textbf{Фазовая коррекция.} Умножение $Y_2[k]$ из~\eqref{eq:y2} на $e^{-j \pi \nu_k}$ для устранения частотно-зависимого сдвига:
    \begin{equation}
    \tilde{Y}_2[k] = Y_2[k] e^{-j \pi \nu_k}.
    \end{equation}
    После компенсации этого сдвига система~\eqref{eq:y1}--\eqref{eq:y2} (при нормировке $\rho_1=\rho_2=0.5$, выносимой в коэффициент $1/\sqrt{2}$) принимает форму матричного умножения на матрицу Адамара размера $2\times2$:

    \begin{equation}
    \begin{pmatrix} Y_1[k] \\ \tilde{Y}_2[k] \end{pmatrix} = 
    \begin{pmatrix} 1 & 1 \\ 1 & -1 \end{pmatrix} 
    \begin{pmatrix} \alpha_1[k] x[k] \\ \alpha_2[k] x[k] \end{pmatrix} + 
    \begin{pmatrix} n_1[k] \\ \tilde{n}_2[k] \end{pmatrix}.
    \end{equation}
    Коэффициенты коррекции $e^{-j \pi \nu_k}$ вычисляются один раз при проектировании и хранятся в ПЗУ; их применение требует всего $O(K)$ комплексных умножений, что составляет менее 5\% от вычислительной стоимости блока БПФ и не вносит существенной задержки.
    
    \item \textbf{Разделение (Обратное преобразование Адамара).} Для сохранения мощности сигнала применяется унитарная форма матрицы Адамара $H_2 = \frac{1}{\sqrt{2}}\begin{bmatrix} 1 & 1 \\ 1 & -1 \end{bmatrix}$. Обратная матрица совпадает с прямой ($H_2^{-1} = H_2^\dagger = H_2$). Использование ортогональных преобразований минимизирует ошибку восстановления сигнала в условиях зашумленного измерения \cite{tropp}. Таким образом:
    \begin{equation}
    \begin{bmatrix} \hat{S}_1[k] \\ \hat{S}_2[k] \end{bmatrix} =\frac{1}{\sqrt{2}}\begin{bmatrix} 1 & 1 \\ 1 & -1 \end{bmatrix}\begin{bmatrix} Y_1[k] \\ \tilde{Y}_2[k] \end{bmatrix},\label{eq:hadamard}
    \end{equation}

    откуда при равномерном распределении мощности ($\rho_1 = \rho_2 = 0.5$):
    \begin{align}
    \hat{S}_1[k] &= \frac{1}{\sqrt{2}} \left( Y_1[k] + \tilde{Y}_2[k] \right) \approx \alpha_1[k] x[k], \\
    \hat{S}_2[k] &= \frac{1}{\sqrt{2}} \left( Y_1[k] - \tilde{Y}_2[k] \right) \approx \alpha_2[k] x[k].
    \end{align}
    Примечание: Множитель $1/\sqrt{2}$ в унитарной матрице $H_2$ точно компенсирует коэффициенты $\sqrt{\rho_1}=\sqrt{\rho_2}=1/\sqrt{2}$ (при активных ВЧ-трактах с $\text{SPLIT}=1$), сохраняя единичную среднюю мощность ветвей и исходную дисперсию шума $\sigma_n^2$.

    Критически важно, что после этого шага шум $n_1$ и $n_2$ остаются независимыми гауссовскими величинами, а сигнальные составляющие $\alpha_1 x$ и $\alpha_2 x$ полностью декоррелированы благодаря ортогональности строк матрицы Адамара. Интерференция между репликами отсутствует.
    
    \item \textbf{Оценка канала (МНК).} В модели используется блочная пилот-структура (пилотные символы $x_p[k]$ на всех $K=64$ поднесущих в начале кадра без частотной интерполяции), откуда получаем оценки $\hat{\alpha}_1[k] = \hat{S}_1[k] / x_p[k]$ и $\hat{\alpha}_2[k] = \hat{S}_2[k] / x_p[k]$. 
    
    \textit{Пояснение:} Оценка канала выполняется \textit{только после} разделения; попытка оценить $\alpha_1$ и $\alpha_2$ до разделения приводит к интерференции пилотов со случайными фазами. Поскольку обе реплики излучаются одним ЦАП и проходят через общий радиотракт, их символьная синхронизация и циклический префикс (CP) совпадают, а остаточная разность задержек/фаз между зонами поглощается комплексными коэффициентами $\alpha_1[k], \alpha_2[k]$ при МНК-оценке \cite{proakis}.

    
    \item \textbf{Комбинирование (MRC).} Итоговая оценка полезного сигнала формируется как взвешенная сумма:
    \begin{equation}
    \hat{x}_{dec}[k] = \frac{\hat{\alpha}_1^*[k] \hat{S}_1[k] + \hat{\alpha}_2^*[k] \hat{S}_2[k]}{|\hat{\alpha}_1[k]|^2 + |\hat{\alpha}_2[k]|^2 + \epsilon},
    \end{equation}
    где $\epsilon$ — малая константа для предотвращения деления на ноль. Эта процедура обеспечивает максимальное отношение сигнал/шум на выходе при независимых замираниях \cite{proakis}.
\end{enumerate}

\subsection{Аналитическая граница производительности}
Для оценки эффективности предложенного алгоритма используется понятие теоретического предела при идеальном знании состояния канала (Perfect CSI). Важно отметить, что данный предел отличается от емкости канала Шеннона, которая определяет максимальную скорость передачи информации. Наш предел отвечает на вопрос: «Какую минимальную вероятность ошибки может обеспечить данная структура приемника (MRC с двумя ветвями), если бы он знал точные значения коэффициентов передачи канала $\alpha_1$ и $\alpha_2$ без ошибок?»

Для рэлеевского канала с порядком разнесения $M=2$ и модуляцией QPSK вероятность битовой ошибки (BER) при использовании MRC с Perfect CSI определяется выражением \cite{proakis}:
\begin{equation}
P_b(\bar{\gamma}) = \left( \frac{1-\mu}{2} \right)^2 (2 + \mu),
\label{eq:ber_theory}
\end{equation}
где $\mu = \sqrt{\frac{\bar{\gamma}}{1+\bar{\gamma}}}$, а $\bar{\gamma}$ — среднее отношение сигнал/шум на одну ветвь разнесения. Эта формула служит ориентиром (кривая ``Идеальный MRC'' в наших экспериментах), к которому аппроксимируются результаты реальной системы, использующей оценку канала по пилотам (LS). Для учёта неидеальности LS-оценки канала с дисперсией ошибки $\sigma_e^2 = c_e \sigma_n^2$ (где $c_e = 1$ при оценке по одиночному пилоту и $c_e = 0.5$ при когерентном усреднении двух пилотных слотов в схеме Аламути) в аналитической модели используется эффективное отношение сигнал/шум на ветвь~\cite{proakis}: \begin{equation} \gamma_{\text{eff}} = \frac{P_{\text{br}}^2}{\sigma_n^2 (P_{\text{br}} + \sigma_e^2) + \kappa P_{\text{br}} \sigma_e^2}, \label{eq:gamma_eff} \end{equation} где $P_{\text{br}}$ --- средняя мощность сигнала на одну ветвь разнесения, $\kappa$ --- число передающих антенн, а слагаемые в знаменателе отражают вклад теплового шума и перекрёстной помехи от ошибки оценки канала. Подстановка $\bar{\gamma} = \gamma_{\text{eff}}$ в~\eqref{eq:ber_theory} (и её обобщение на произвольный порядок $L$) задаёт аналитический ориентир с учётом потерь LS-оценки.

\section{Методология исследования}

Верификация архитектуры выполнена методом Монте-Карло на символьном уровне (после операции БПФ в приёмнике OFDM). Модель учитывает:
\begin{itemize}
    \item \textbf{Модуляция:} QPSK на каждой из $K=64$ поднесущих OFDM, $E|x[k]|^2=1$.
    \item \textbf{Канал:} Независимые рэлеевские замирания, моделируемые как циркулярно-симметричные комплексные гауссовские случайные величины $\alpha_i[k] \sim \mathcal{CN}(0, 1)$ ($\E{|\alpha_i[k]|^2}=1$), где $i \in \{1, 2\}$ --- индекс спектральной реплики (зоны Найквиста), а $k \in \{0, \dots, K-1\}$ --- индекс поднесущей.
    \item \textbf{Шум:} Аддитивный белый гауссов шум с дисперсией $\sigma_n^2$, определяемой заданным SNR ($SNR = E_s/\sigma_n^2$).
    \item \textbf{Нормировка мощности:} Полная передаваемая мощность нормирована к единице ($\rho_1 + \rho_2 = 1$), распределение между репликами равномерное ($\rho_1=\rho_2=0.5$).
    \item \textbf{Объем выборки:} Адаптивное накопление до достижения заданного числа эффективных отсчётов ($N_{\text{eff}} \ge 400$): от $4\cdot10^3$ до $6\cdot10^5$ OFDM-блоков (от $2.56\cdot10^5$ до $3.84\cdot10^7$ символов QPSK) на точку SNR с отбраковкой точек при $N_{\text{eff}} < 150$.
    \item \textbf{Диапазон SNR:} $0$--$40$~дБ с шагом $2$~дБ (табличный анализ в рабочей области $8$--$24$~дБ).

\end{itemize}

\subsection{Обоснование выбора базовых линий сравнения}
Для объективной оценки вклада именно алгоритма разнесения необходимо изолировать его эффект от влияния аппаратурных затрат (наличия второго АЦП). Поэтому выбраны следующие линии сравнения:

\begin{enumerate}
    \item \textbf{Базовая схема (однозонный прием, 1 АЦП)}: Однозонный OFDM, 1 АЦП. Порядок разнесения $d \approx 1$. Является эталоном простейшего приемника, не имеющего дополнительных аппаратных средств.
    \item \textbf{Контрольная схема (однозонный прием, 2 АЦП)}: Однозонный OFDM, 2 АЦП со сдвигом $T_s/2$ и когерентным усреднением одного и того же сигнала. Эта схема имеет то же количество АЦП, что и предлагаемая, но \textit{не использует} спектральные реплики для создания разнообразия каналов. Она изолирует процессный выигрыш от наличия второго АЦП (снижение шума в $\sqrt{2}$ раз), сохраняя порядок разнесения $d \approx 1$ \cite{walden}. Сравнение с этой линией критически важно, чтобы исключить обвинение в том, что выигрыш получен просто за счет удвоения количества аппаратуры.
    \item \textbf{Предлагаемая архитектура (ПСИР)}: Двухрепликовый OFDM, 2 АЦП, разделение Адамаром, оценка канала методом наименьших квадратов, MRC.
    \item \textbf{Предлагаемая архитектура с неидеальностями}: То же, что и предлагаемая архитектура, но с учётом неидеальностей TI-ADC: рассогласование усиления $2\%$ и фазовая ошибка тактирования $2.3^\circ$ (что эквивалентно суммарному перекосу выборки и апертурному джиттеру $\delta t \approx 6.4{\cdot}10^{-3} T_s$, превышающему типовые допуски современных микросхем TI-ADC \cite{ieee_std,adi_mt008,razavi}), без калибровки матрицы разделения. Воспроизводимость всех случайных реализаций обеспечена фиксацией начального состояния генератора псевдослучайных чисел (\texttt{seed}) в коде симуляции \cite{github}.
    \item \textbf{Идеальный MRC (предел при совершенном CSI)}: Верхняя граница производительности. Предполагается идеальное знание канала $\alpha_1, \alpha_2$ при идеальном разделении реплик. Соответствует формуле \eqref{eq:ber_theory} \cite{proakis}.
\end{enumerate}

\textbf{Проверка корректности модели.} На этапе предварительной верификации модели выполнена проверка, что фазовое соотношение $\pi$ верно не только для комплексной базовой модели, но и для вещественного широкополосного сигнала после выделения положительной частотной компоненты. Для сигнала \eqref{eq:real_signal} с восемью тонами, дискретизированного со сдвигом $T_s/2$, восстановление $\alpha_1,\alpha_2$ матрицей Адамара дало максимальную ошибку $|\varepsilon_{\alpha_1}| = 1{,}31\cdot10^{-13}$, $|\varepsilon_{\alpha_2}| = 4{,}49\cdot10^{-13}$, что подтверждает корректность комплексной модели как комплексного эквивалента вещественного тракта. Подробности эксперимента, включая исходный код, приведены в репозитории~\cite{github}.

\subsection{Машинная верификация ключевых утверждений}
Для обеспечения строгой воспроизводимости и исключения скрытых допущений в симулятор встроены девять автоматических тестов (\texttt{passband}, \texttt{power}, \texttt{pipeline}, \texttt{mismatch}, \texttt{models}, \texttt{ls\_theory}, \texttt{estimator}, \texttt{asymptote}, \texttt{error\_floor}), которые выполняются \emph{до} основного прогона и останавливают его при провале любого из условий:

\begin{enumerate}
\item \textbf{Вещественный тракт (\texttt{passband}) и бюджет мощности (\texttt{power}).} Подтверждено совпадение вещественного полосового тракта~\eqref{eq:real_signal} с комплексной базовой моделью с точностью до $4.49\cdot10^{-13}$, а также выполнение всех шести тождеств энергетического баланса и сохранения дисперсии шума $\sigma_n^2$ после унитарного преобразования Адамара~\eqref{eq:hadamard} с точностью $\le 3\%$. При активных ВЧ-трактах ($\text{SPLIT}=1$) каждый АЦП несёт полную мощность сигнала, тогда как при пассивном делителе ($\text{SPLIT}=1/\sqrt{2}$) выигрыш по отношению к классическому ЧР на 4 АЦП обнуляется ($0.00$~дБ).

\item \textbf{Порядок оценки канала (\texttt{pipeline}).} Сравниваются два конвейера при $\SNR=10$~дБ: корректный (разделение Адамара $\to$ LS по поднесущим) и наивный (оценка смешанных наблюдений АЦП до разделения). Корректный конвейер достигает теоретического предела $\text{MSE} = 0.1000 = \sigma_n^2$ и $\BER = 1.828\cdot10^{-2}$, тогда как наивный деградирует до $\text{MSE} = 0.6853$ и $\BER = 4.998\cdot10^{-1}$ (ухудшение в $\mathbf{27.3\times}$). Это строго доказывает, что разделение должно предшествовать оценке канала.

\item \textbf{Корреляция и отсутствие порога ошибок при рассогласовании TI-ADC (\texttt{mismatch}, \texttt{error\_floor}).} Измеренная остаточная корреляция ветвей и подавление перекрёстных помех ($\text{XT} = 33.00$~дБ при рассогласовании $2\%$ и $2.3^\circ$) точно совпадают с аналитикой. Проверка на $1.2\cdot10^5$ блоках при $\SNR = 30$~дБ подтверждает отсутствие порога ошибок (\emph{error floor}): $\BER_{\text{ideal}} = 2.853\cdot10^{-6}$ против $\BER_{\text{mismatch}} = 2.859\cdot10^{-6}$ ($\Delta = +0.23\%$).

\item \textbf{Согласованность аналитики и асимптотик (\texttt{models}, \texttt{ls\_theory}, \texttt{estimator}, \texttt{asymptote}).} Верифицированы нулевая ошибка при отсутствии шума для всех 10 схем, совпадение полуаналитической оценки (условная BER) с жёсткими решениями, точность аналитической модели~\eqref{eq:ber_theory} с учётом потерь LS-оценки $\gamma_{\text{eff}}$~\eqref{eq:gamma_eff} и выход асимптотических наклонов в окне $60$--$100$~дБ строго на конструктивный порядок $L \in \{1, 2, 4, 8\}$ с погрешностью $<0.03$ (рис.~\ref{fig:slope_ladder}).

\item \textbf{Зависимость порядка от независимости реплик.} Свип частотной корреляции реплик $\rho_f$ в гибриде $\text{Alamouti}+\text{ПСИР}$ показывает, что при $\rho_f = 0$ локальный наклон в окне $12$--$20$~дБ составляет $d_{\text{eff}} = 7.22$, а при росте корреляции до $\rho_f = 0.99$ снижается до $d_{\text{eff}} = 3.94$ ($L=8 \to L=4$), подтверждая физическую природу удвоения порядка.
\end{enumerate}

Полный протокол всех проверок с конфигурацией, статусами \texttt{ok/FAIL}, результатами идентификации $\hat{L}$ и таблицей наклонов экспортируется в файл \texttt{protocol.txt}.


\section{Результаты символьной верификации}

\subsection{Порядок разнесения и BER}
На рисунке~\ref{fig:ber_curves} представлены кривые вероятности битовой ошибки (BER) для рассмотренных архитектур. Таблица~\ref{tab:results} содержит сводные метрики, полученные на основе финального запуска симуляции.

\begin{figure}[H]
    \centering
    \includegraphics[width=0.95\textwidth, keepaspectratio]{psir_ber.png} 
    \caption{Зависимость BER от SNR для различных архитектур приёма (линии --- точная теория, маркеры --- измеренные надёжные точки с 95\% ДИ)}
    \label{fig:ber_curves}
\end{figure}

Как видно из рисунка~\ref{fig:ber_curves} и таблицы~\ref{tab:results}, кривая предлагаемой архитектуры ПСИР идентифицируется по форме всей кривой как система с порядком разнесения $\hat{L}=2$ (локальный наклон $d = 1.872$ в окне 14--24~дБ и $d = 2.102$ в окне 22--30~дБ), в то время как базовая ($d=0.968$) и контрольная ($d=0.982$) схемы демонстрируют порядок $\hat{L}=1$. Неидеальная схема с рассогласованием TI-ADC 2\%/2.3$^\circ$ ($d=1.871$) и классическое ЧР с 4 АЦП ($d=1.916$) практически накладываются на кривую ПСИР: при сравнении на общих случайных числах (CRN) максимальное отклонение $|\Delta\BER/\BER|$ не превышает $0.25\%$ во всём рабочем диапазоне 0--30~дБ (рис.~\ref{fig:robust}).

\begin{figure}[H]
    \centering
    \includegraphics[width=0.9\textwidth, keepaspectratio]{psir_robust.png}
    \caption{Робастность ПСИР к рассогласованию усиления ($2\%$) и фазовой ошибке тактирования ($2.3^\circ$) TI-ADC на общих случайных числах (CRN): кривые BER (сверху) и относительная разность $\Delta\BER/\BER$ в процентах с $95\%$ доверительным интервалом (снизу)}
    \label{fig:robust}
\end{figure}


При этом важно отметить, что кривые ``Неидеальная'' (с рассогласованием 2\% и 2.3°) и ``Классическое частотное разнесение'' визуально совпадают с кривой ``Предлагаемая архитектура'' в рабочем диапазоне SNR 6--24 дБ. Это совпадение имеет два важных следствия: во-первых, оно подтверждает математическую эквивалентность ПСИР классическому разнесению благодаря унитарности матрицы Адамара; во-вторых, оно доказывает робастность метода, так как перекрестные помехи от неидеальностей АЦП находятся ниже уровня теплового шума и не формируют порога ошибок (error floor) в рабочей зоне.

\begin{table}[htbp]
\centering
\footnotesize
\caption{Сводные характеристики и энергетические выигрыши исследуемых архитектур (QPSK, $K{=}64$, независимые рэлеевские замирания по антеннам и зонам Найквиста, $\rho_1{=}\rho_2{=}0.5$)}


\label{tab:results}
\begin{adjustbox}{max width=\textwidth}
\begin{tabular}{@{}lccccccccc@{}}
\toprule
\multicolumn{10}{l}{\textbf{Часть A. SISO-схемы и эталоны}} \\
\midrule
\textbf{Схема} & \textbf{АЦП} & \textbf{$L$ / $\hat{L}$} & \textbf{SE} & \textbf{$d_{\text{low}}$ (окно)} & \textbf{$d_{\text{high}}$ (окно)} & \multicolumn{3}{c}{\textbf{BER при SNR, дБ}} & \textbf{Выигрыш vs контроль} \\
\cmidrule(lr){7-9}
 & & & & & & 14 & 18 & 22 & @ $10^{-2}$ / @ $10^{-4}$ (LS), дБ \\
\midrule
OFDM базовый (1 зона) & 1 & 1 / 1 & 1.00 & 0.968 (14--24) & 0.976 (24--34) & $3.63{\cdot}10^{-2}$ & $1.54{\cdot}10^{-2}$ & $6.31{\cdot}10^{-3}$ & --- \\
OFDM + TI-усреднение (контроль) & 2 & 1 / 1 & 1.00 & 0.982 (14--24) & 0.994 (24--34) & $1.91{\cdot}10^{-2}$ & $7.63{\cdot}10^{-3}$ & $3.10{\cdot}10^{-3}$ & --- \\
\textbf{ПСИР (предлагаемая)} & \textbf{2} & \textbf{2 / 2} & \textbf{0.50} & \textbf{1.872 (14--24)} & \textbf{2.102 (22--30)} & $\mathbf{3.81{\cdot}10^{-3}}$ & $\mathbf{6.82{\cdot}10^{-4}}$ & $\mathbf{1.17{\cdot}10^{-4}}$ & \textbf{5.24 / 14.67} \\
ПСИР + неидеальности (2\%, 2.3$^\circ$) & 2 & 2 / 2 & 0.50 & 1.871 (14--24) & 2.100 (22--30) & $3.81{\cdot}10^{-3}$ & $6.83{\cdot}10^{-4}$ & $1.17{\cdot}10^{-4}$ & 5.24 / 14.67 \\
Классич. ЧР (2 АЦП) & 2 & 2 / 2 & 0.50 & 1.847 (14--24) & 1.984 (22--30) & $1.29{\cdot}10^{-2}$ & $2.52{\cdot}10^{-3}$ & $4.61{\cdot}10^{-4}$ & 2.21 / 11.66 \\
Классич. ЧР (4 АЦП) & 4 & 2 / 2 & 0.50 & 1.916 (14--24) & 2.027 (22--30) & $3.79{\cdot}10^{-3}$ & $6.85{\cdot}10^{-4}$ & $1.16{\cdot}10^{-4}$ & 5.29 / 14.67 \\
ПСИР, идеальное CSI & 2 & 2 / 2 & 0.50 & 1.924 (14--24) & 2.014 (22--30) & $1.02{\cdot}10^{-3}$ & $1.82{\cdot}10^{-4}$ & $2.82{\cdot}10^{-5}$ & 8.25 / 17.70 \\
\midrule
\multicolumn{10}{l}{\textbf{Часть B. Гибриды с пространственным разнесением}} \\
\midrule
\textbf{Схема} & \textbf{Конфиг.} & \textbf{АЦП} & \textbf{$L$ / $\hat{L}$} & \textbf{SE} & \textbf{$d_{\text{low}}$ (окно)} & \textbf{$d_{\text{high}}$ (окно)} & \textbf{BER@14 дБ} & \multicolumn{2}{c}{\textbf{Выигрыш vs ПСИР (LS), дБ}} \\
\cmidrule(lr){9-10}
 & & & & & & & & @ $10^{-3}$ & @ $10^{-4}$ / @ $10^{-6}$ \\
\midrule
\textbf{ПСИР + SIMO $1{\times}2$} & $\mathbf{1{\times}2}$ & \textbf{4} & \textbf{4 / 4} & \textbf{0.50} & \textbf{3.011 (6--14)} & \textbf{3.825 (14--22)} & $\mathbf{5.89{\cdot}10^{-5}}$ & \textbf{6.90} & \textbf{9.04 / 13.70} \\
Alamouti $2{\times}2$ без ПСИР & $2{\times}2$ & 4 & 4 / 4 & 1.00 & 2.430 (2--12) & 3.594 (12--20) & $5.30{\cdot}10^{-5}$ & --- & --- \\
\textbf{Alamouti $2{\times}2$ + ПСИР} & $\mathbf{2{\times}2}$ & \textbf{4} & \textbf{8 / 8} & \textbf{0.50} & \textbf{3.477 (2--12)} & \textbf{7.423* (16--24)} & $\mathbf{1.40{\cdot}10^{-6}}$ & \textbf{8.55} & \textbf{11.54 / 18.12} \\
\midrule
\multicolumn{10}{l}{\hspace{0.5em}$\uparrow$ \textit{При равных 4 АЦП: Alamouti даёт $L{=}4$, ПСИР даёт $L{=}2$, их комбинация $\Rightarrow$ $L{=}4{\times}2{=}8$ (мультипликативность ценой $\text{SE}/2$)}} \\
\midrule
\multicolumn{10}{l}{\textbf{Часть C. Энергетический выигрыш ПСИР при заданном уровне BER, дБ}} \\
\midrule
\multicolumn{5}{l}{\textbf{Пара сравнения}} & \multicolumn{2}{c}{\textbf{Измерено @ BER}} & \multicolumn{3}{c}{\textbf{Теория LS (идеал. CSI) @ BER}} \\
\cmidrule(lr){6-7} \cmidrule(lr){8-10}
\multicolumn{5}{l}{} & $10^{-2}$ & $10^{-4}$ & $10^{-3}$ & $10^{-4}$ & $10^{-6}$ \\
\midrule
\multicolumn{5}{l}{ПСИР vs Контроль аппаратуры (OFDM + TI-усреднение, 2 АЦП)} & 5.24 & 14.56 & 9.82 (9.86) & 14.67 (14.69) & 24.61 (24.61) \\
\multicolumn{5}{l}{ПСИР vs OFDM базовый (1 зона, 1 АЦП)} & 8.25 & --- & 12.83 (12.87) & 17.68 (17.70) & 27.62 (27.62) \\
\multicolumn{5}{l}{ПСИР vs Классич. ЧР (2 АЦП)} & 3.03 & 2.82 & 3.01 (3.01) & 3.01 (3.01) & 3.01 (3.01) \\
\multicolumn{5}{l}{ПСИР vs Классич. ЧР (4 АЦП, аппаратный эталон)} & $-0.05$ & $-0.07$ & 0.00 (0.00) & 0.00 (0.00) & 0.00 (0.00) \\
\multicolumn{5}{l}{\textbf{Alamouti $2{\times}2$ + ПСИР vs Alamouti $2{\times}2$ без ПСИР (равные 4 АЦП)}} & \textbf{0.70} & \textbf{2.43} & \textbf{1.55 (1.73)} & \textbf{2.45 (2.57)} & \textbf{4.41 (4.48)} \\
\multicolumn{5}{l}{Alamouti $2{\times}2$ + ПСИР vs ПСИР SISO} & 5.87 & 11.64 & 8.55 (8.79) & 11.54 (11.70) & 18.12 (18.20) \\
\bottomrule
\end{tabular}
\end{adjustbox}
\end{table}

\textbf{Обозначения и пояснения к табл.~\ref{tab:results}.} В таблице приняты следующие обозначения: $L$ --- конструктивный порядок разнесения (число независимых ветвей); $\hat{L}$ --- порядок, идентифицированный по форме всей кривой Монте-Карло; $\text{SE}$ --- спектральная эффективность (бит/с/Гц); $d_{\text{low}}$ и $d_{\text{high}}$ --- локальные наклоны $\lg\BER$ от $\lg\SNR$ в низком и высоком окнах SNR (звёздочкой~* отмечен аналитический наклон по модели~\eqref{eq:gamma_eff} за пределами границы надёжности Монте-Карло). Сравнение ключевых архитектур выполняется при строго одинаковом количестве АЦП (2 АЦП для SISO и 4 АЦП для MIMO~$2{\times}2$), что изолирует эффект ПСИР от аппаратного выигрыша за счёт удвоения числа конвертеров. Расхождение локальных наклонов $d_{\text{low}}, d_{\text{high}}$ с конструктивным порядком $L$ (особенно при $L=4, 8$) объясняется тем, что на рабочих SNR кривая MRC находится на переходном участке и выходит на асимптоту $d \to L$ только при $\BER \to 0$; поэтому достоверность порядка подтверждается совпадением $\hat{L} = L$ по форме всей кривой. В части~C значения в скобках соответствуют пределу идеального знания канала (Perfect CSI), а значения перед скобками --- аналитической модели~\eqref{eq:gamma_eff} с учётом шума LS-оценки.

\textbf{Методика извлечения порядка разнесения.} 

Порядок частотного разнесения $d$ определялся как асимптотический наклон 
зависимости вероятности битовой ошибки от отношения сигнал/шум в логарифмических 
координатах. В соответствии с классическим высокосигнальным асимптотическим выражением для 
рэлеевского канала с разнесением \cite{proakis,goldsmith}:
\begin{equation}
P_b(\bar{\gamma}) \approx C \cdot \bar{\gamma}^{-d},
\label{eq:diversity_asymptotic}
\end{equation}
где $P_b$ --- вероятность битовой ошибки, $\bar{\gamma}$ --- номинальное среднее отношение сигнал/шум 
на ветвь в линейном масштабе ($\bar{\gamma} \gg 1$), а $C$ --- константа (коэффициент энергетического сдвига), зависящая от типа модуляции, порядка разнесения $L$ и точности оценки канала: при идеальном CSI $C = \binom{2L-1}{L} 2^{-L}$ для QPSK~\cite{proakis,goldsmith}, а при неидеальной LS-оценке~\eqref{eq:gamma_eff} величина $C$ возрастает в $(1 + \kappa c_e)^L$ раз (где $\kappa$ --- число передающих антенн, $c_e$ --- нормированная дисперсия ошибки оценки канала).

Логарифмирование обеих частей по основанию 10 дает линейную зависимость:
\begin{equation}
\log_{10}(P_b) \approx \log_{10}(C) - d \cdot \log_{10}(\bar{\gamma}).
\label{eq:diversity_log}
\end{equation}
Аппроксимация методом наименьших квадратов (линейная регрессия) выполнялась по надёжным точкам Монте-Карло ($N_{\text{eff}} \ge 150$) в низких и высоких окнах SNR (табл.~\ref{tab:results}), а конструктивный порядок $\hat{L}$ дополнительно верифицировался по минимальному среднеквадратичному отклонению всей кривой от аналитической модели MRC. Доверительные интервалы для параметра $d$ рассчитывались как $\pm 1.96 \cdot \sigma_{slope}$ на уровне доверия 95\%. Эволюция локального наклона $d(\SNR) \to L$ с ростом отношения сигнал/шум проиллюстрирована на рисунке~\ref{fig:slope_ladder}.

\begin{figure}[H]
    \centering
    \includegraphics[width=0.9\textwidth, keepaspectratio]{psir_slope.png}
    \caption{Рост локального наклона $d = -\mathrm{d}\lg\BER/\mathrm{d}\lg\SNR$ к конструктивному порядку $L$ с увеличением SNR}
    \label{fig:slope_ladder}
\end{figure}

\noindent\textbf{Пояснение к рис.~\ref{fig:slope_ladder}.} Сплошные и штриховые кривые представляют собой точный аналитический наклон $d(\SNR) = -\mathrm{d}\lg P_b / \mathrm{d}\lg\bar{\gamma}$ для когерентного комбинирования (MRC) в рэлеевском канале~\eqref{eq:ber_theory}. Он вычислен путём дифференцирования теоретической зависимости BER с подстановкой эффективного отношения сигнал/шум $\gamma_{\text{eff}}$~\eqref{eq:gamma_eff}, учитывающего ошибку МНК-оценки канала по пилотам~\cite{proakis}. Маркеры с доверительными интервалами ($\pm 1.96\,\sigma_{\text{slope}}$, уровень доверия $95\%$) отражают экспериментальный наклон, полученный методом линейной регрессии в логарифмических координатах~\eqref{eq:diversity_log} по точкам моделирования Монте-Карло с шагом $\Delta\SNR = 2$~дБ. Небольшое смещение маркеров ниже аналитических кривых обусловлено двумя факторами: во-первых, регрессия по дискретным узлам с шагом $2$~дБ даёт среднюю крутизну на конечном отрезке, которая на участке роста систематически ниже мгновенной производной; во-вторых, у границы измеримости (для $\text{Alamouti} + \text{ПСИР}$ при $\SNR = 14$~дБ, где $\BER \approx 1.4{\cdot}10^{-6}$) сказывается статистическая флуктуация выборки у порога надёжности $N_{\text{eff}} \ge 150$.

\subsection{Интерпретация результатов}

Для практического инженера-связиста ключевыми являются не столько параметры регрессии ($d, R^2$), сколько конкретные значения снижения требуемого отношения сигнал/шум (выигрыш по энергии).

\textbf{1. Энергетический бюджет и дальность связи.}
Как следует из таблицы~\ref{tab:results}, при целевом уровне надежности $\BER=10^{-2}$ (типично для незащищенных голосовых или телеметрических каналов IoT) предлагаемая архитектура требует на \textbf{5.24 дБ} (5.26~дБ по теории LS) меньше мощности передатчика по сравнению с контрольной схемой (два АЦП, но без использования реплик).
По сравнению с самой простой базовой схемой (один АЦП), измеренный выигрыш составляет \textbf{8.25 дБ} (8.27~дБ по теории LS).

В линейном масштабе (в разах) это означает, что для достижения той же вероятности ошибки \textbf{мощность} передатчика можно снизить примерно в \textbf{6.7 раза} ($10^{8.25/10} \approx 6.68$), а эквивалентная \textbf{дальность} связи при фиксированной мощности возрастает примерно в \textbf{2.6 раза} ($10^{8.25/20} \approx 2.59$). С учётом потерь в кабеле и КПД усилителя мощности реальная экономия батареи окажется ближе к 3--4-кратной. Для автономных датчиков это прямое увеличение срока службы аккумуляторов или расширение зоны покрытия без установки повторителей.

\textbf{2. Качество связи и отсутствие интерференции.}
Результаты показывают, что кривая предлагаемой архитектуры практически сливается с кривой неидеальной схемы ($d_{\text{low}}=1.872$ против $1.871$ и $d_{\text{high}}=2.102$ против $2.100$, макс. $|\Delta\BER/\BER| = 0.25\%$ на CRN, рис.~\ref{fig:robust}). Это доказывает, что метод устойчив к реальным технологическим допускам. Более того, вертикальный сдвиг между предлагаемой архитектурой (LS-оценка~\eqref{eq:gamma_eff}) и теоретическим пределом идеального MRC~\eqref{eq:ber_theory} (genie-режим, рис.~\ref{fig:ber_curves}) составляет около 2--3~дБ в рабочей области SNR и монотонно уменьшается с ростом SNR по мере того, как шум оценки канала становится пренебрежимой долей суммарного шума. Этот небольшой зазор обусловлен исключительно шумом оценки канала по пилотам, а не остаточной интерференцией между репликами. Если бы матрица Адамара не обеспечивала полную ортогональность, мы наблюдали бы ошибочный порог (error floor) на высоких SNR или значительное расхождение с Идеальным MRC даже при малых неидеальностях. Отсутствие такого порога подтверждает корректность физической модели разделения.

\textbf{3. Мультипликативный эффект с пространственным разнесением.}
Предлагаемый метод комплементарен классическим пространственным схемам при одинаковом аппаратном бюджете (4 АЦП: по два интерливинговых АЦП со сдвигом $T_s/2$ на каждую из двух приёмных антенн). В многоантенной конфигурации каналы между всеми парами передающих и приёмных антенн моделируются как независимые рэлеевские ($\mathcal{CN}(0,1)$); алгоритм ПСИР (фазовая коррекция и преобразование Адамара $H_2$) применяется независимо на выходе каждой приёмной антенны до пространственно-временного декодирования Аламути и MRC, а раздельная МНК-оценка коэффициентов от двух передающих антенн выполняется по паре ортогональных пилотных символов $[1, 1]^T$ и $[1, -1]^T$~\cite{proakis}. Численно верифицированы три гибридных конфигурации:

\begin{itemize}
\item \textbf{ПСИР + SIMO $1{\times}2$} (4 АЦП): идентификация по кривой даёт $\hat{L}=4$, локальный наклон растёт от $d_{\text{low}} = 3.011$ (окно 6--14~дБ) до $d_{\text{high}} = 3.825$ (окно 14--22~дБ), асимптота $L=4$.
\item \textbf{MIMO $2{\times}2$ Alamouti без ПСИР} (4 АЦП): $\hat{L}=4$, локальный наклон растёт от $d_{\text{low}} = 2.430$ (окно 2--12~дБ) до $d_{\text{high}} = 3.594$ (окно 12--20~дБ), асимптота $L=4$.
\item \textbf{MIMO $2{\times}2$ Alamouti + ПСИР} (4 АЦП): $\hat{L}=8$, локальный наклон растёт от $d_{\text{low}} = 3.477$ (окно 2--12~дБ) до $d_{\text{high}} = 7.423$ (окно 16--24~дБ), асимптота $L=8$.
\end{itemize}

На крутом участке 16--24~дБ добавление ПСИР к схеме Alamouti при тех же 4 АЦП удваивает крутизну кривой ($3.86 \to 7.42$, отношение $1.92\times$), а выигрыш по SNR при $\BER=10^{-4}$ составляет $\mathbf{2.43}$~дБ (измеренный) и $\mathbf{2.45}$~дБ (по теории LS), что строго подтверждает мультипликативность разнесений ($L = 4 \times 2 = 8$, рис.~\ref{fig:mult}). При росте частотной корреляции реплик $\rho_f \to 0.99$ эффективный наклон гибрида падает с $7.22$ до $3.94$, доказывая, что удвоение порядка обусловлено именно статистической независимостью зон Найквиста. Практическая реализация подобных гибридных схем обсуждается в работах по субнайквистовой радарной съемке и аппаратной реализации на FPGA \cite{han,gadawe}.

\begin{figure}[H]
    \centering
    \includegraphics[width=0.88\textwidth, keepaspectratio]{psir_mult.png}
    \caption{Мультипликативное удвоение порядка разнесения при равных 4 АЦП: переход от Alamouti $2{\times}2$ ($L=4$) к Alamouti $2{\times}2$ + ПСИР ($L=8$)}
    \label{fig:mult}
\end{figure}

\section{Обсуждение и ограничения}

\textbf{Инновационный потенциал и экономика системы.}
Основная ценность предложенного метода заключается не в перераспределении вычислительной сложности, а в монетизации физического артефакта дискретизации. Традиционно спектральные копии считаются потерями, требующими дорогостоящих аналоговых фильтров для их удаления. Предлагаемый подход позволяет использовать более простые аналоговые фильтры на стороне передатчика, однако требует наличия диплексера или полосового фильтра для подавления зеркального образа ($f_s-f_0$) и интермодуляционных искажений, которые не могут быть устранены цифровым разделением. Это приводит к снижению общего веса, габаритов и энергопотребления системы связи, особенно в сегменте IoT и мобильной связи, где размер батареи и антенны являются лимитирующими факторами.

\textbf{Синергия с пространственным разнесением (MIMO).}
Предложенный метод не конкурирует с технологией MIMO, а дополняет ее, используя ортогональные ресурсы. MIMO обеспечивает разнесение за счет пространственного разнообразия (различные пути распространения волн к разным антеннам), тогда как наш метод обеспечивает разнесение за счет частотного разнообразия (различные участки спектра радиолинии, разделенные на расстояние $f_s$). Поскольку эти механизмы независимы, их ветви разнесения объединяются мультипликативно по числу степеней свободы ($L_{\text{total}} = L_{\text{space}} \cdot L_{\text{freq}}$). В системе $2 \times 2$ MIMO с применением предлагаемого алгоритма число независимых ветвей разнесения увеличивается с 4 (чистый Alamouti) до 8 (Alamouti + ПСИР), что повышает конструктивный порядок разнесения с $L=4$ до $L=8$. Измеренный локальный наклон в окне 2--12~дБ составляет $d=3.477$ для гибрида и $d=2.430$ для чистого Alamouti (на крутом участке 16--24~дБ --- $7.423$ против $3.856$, рис.~\ref{fig:slope_ladder}), а разность SNR при $\BER=10^{-4}$ равна $\mathbf{2.43}$~дБ (2.45~дБ по теории LS, табл.~\ref{tab:results} и рис.~\ref{fig:mult}). Это позволяет достигать уровней надежности, характерных для крупногабаритных многоантенных массивов, используя компактные однопоточные или малоантенные приемники, что критически важно для бортовой аппаратуры дронов, спутников LoRa/IoT и мобильных терминалов специального назначения.

\textbf{Управление занятой полосой пропускания.}
В базовой реализации, представленной в данной работе, спектральные реплики размещаются в удаленных зонах Найквиста (разнос $f_s$), что обеспечивает максимальную статистическую независимость каналов ($f_s > B_c$), но приводит к расширению занятой полосы. Однако архитектура ПСИР допускает модификацию передатчика для уплотнения спектра. Путем разбиения информационной нагрузки на два параллельных потока, размещаемых в смежных зонах Найквиста (например, с частотным разносом, равным ширине полосы одного потока), можно сохранить общую ширину канала неизменной. В этом режиме метод переходит из категории «Frequency Diversity» (повышение надежности ценой полосы) в категорию «Hybrid Multiplexing-Diversity», обеспечивая баланс между скоростью и помехоустойчивостью. Математика разделения (матрица Адамара и фазовая коррекция) остается валидной, так как она зависит исключительно от детерминированного сдвига тактов $\Delta t = T_s/2$, а не от абсолютного расстояния между зонами Найквиста в эфире. Таким образом, проблема «пустых промежутков» в спектре решается тривиальным изменением стратегии мультиплексирования на передатчике без усложнения приемного тракта.

\textbf{Компромисс спектральной эффективности.} Метод требует передачи дублированной информации в двух зонах Найквиста, что снижает эффективную спектральную эффективность (бит/с/Гц) примерно в 2 раза. Поэтому предлагаемая архитектура не рекомендуется для коммерческих сетей мобильной связи, оптимизированных на максимум пропускной способности. Область, где метод полезен — системы, где спектр относительно дешев или избыточен (ISM диапазоны, спутниковые транспондеры, выделенные каналы спецсвязи), но жестко ограничены масса, габариты и энергопотребление приемника. В этих случаях выигрыш в надежности (diversity gain) и упрощение конструкции (отсутствие второго ВЧ-тракта) перевешивают потери в скорости передачи данных.

\textbf{Масштабируемость архитектуры.} Хотя в данной работе верифицирован базовый случай $M=2$ (две зоны Найквиста), предложенный принцип детерминированного разделения реплик масштабируется на произвольное число зон $M$ путем замены матрицы Адамара на полифазное дискретное преобразование Фурье (DFT) или обобщенные ортогональные матрицы. При увеличении $M$ растут требования к точности синхронизации тактовых сигналов ($\delta t < T_s/(4M)$), однако современные технологии распределения clock-сигнала в FPGA/ASIC позволяют реализовать такие схемы для $M \le 4$ без существенного ухудшения характеристик. Расширение числа зон Найквиста позволяет получить порядок разнесения $d \approx M$ (по аналогии с многопорядковым частотным разнесением в OFDM \cite{abdelgader}, но за счёт реплик дискретизации вместо выделения отдельных поднесущих), что делает метод конкурентоспособным с многоантенными системами в условиях ограниченного физического пространства.

\textbf{PAPR и динамический диапазон.} На передающей стороне спектральная реплика является естественным образом дискретизации ЦАП, поэтому пик-фактор (PAPR) непосредственно на выходе ЦАП не увеличивается; однако при совместном усилении обеих зон в широкополосном усилителе мощности (УМ) и при наложении реплик на входе АЦП возрастают требования к линейности и динамическому диапазону тракта (back-off УМ и разрядности АЦП). Детальный анализ нелинейных искажений и PAPR во временной области при масштабировании числа поднесущих ($K = 256\text{--}4096$) и переходе к плотным созвездиям (64-QAM) выходит за рамки символьной модели и составляет предмет дальнейших исследований. Физические ограничения аналоговых узлов (апертурный джиттер, тепловые шумы) подробно описаны в монографии Разаи~\cite{razavi}; они снижают эффективное SNR, но не нарушают ортогональность разделения реплик в цифровой области при стабильной тактовой синхронизации.

\textbf{Условия применимости.} Метод эффективен только при выполнении условия $f_s > B_c$, где $B_c \approx (5\tau_{\text{rms}})^{-1}$ --- когерентная ширина полосы канала по уровню корреляции $0.5$ \cite{goldsmith}. При слабой многолучевости ($\tau_{\text{rms}} \to 0$) реплики становятся коррелированными ($\rho_f \to 1$), и выигрыш частотного разнесения вырождается (что подтверждено тестом свипа $\rho_f$ в разд.~3.2, где порядок гибрида снижается с $L=8$ до $L=4$). В типичных наземных многолучевых каналах (городская застройка, помещения, пересечённая местность) разброс задержек составляет $\tau_{\text{rms}} \sim 0.1\text{--}1~\text{мкс}$, чему соответствует $B_c \sim 0.2\text{--}2~\text{МГц}$; следовательно, уже при частоте дискретизации $f_s \ge 10~\text{МГц}$ разнос реплик с запасом превышает $B_c$. Напротив, для трасс с прямой видимостью и малым разбросом задержек ($\tau_{\text{rms}} \lesssim 20~\text{нс}$, $B_c \gtrsim 10~\text{МГц}$) требуется пропорциональное повышение $f_s$.

\textbf{Занятая полоса и регуляторные маски.} Использование второй реплики увеличивает занятую полосу примерно до $f_s+B$; необходимо проверять соответствие регуляторным маскам внеполосных излучений. Аналоговая полоса пропускания АЦП, ЦАП, усилителей и фильтров должна обеспечивать линейную работу в полосе реплик, например $[f_s, f_s+B]$ для $M=2$. Внеполосные помехи и блокирующие сигналы, попадающие в полосы алиасинга, не разделяются матрицей Адамара и требуют достаточной селективности фронтенда и высокого динамического диапазона АЦП.

\section{Заключение}

В статье предложена и численно верифицирована на символьном уровне сквозная архитектура передачи и приема OFDM-сигналов с использованием детерминированного частотного разнесения на основе спектральных образов ЦАП и интерливинговых АЦП (метод ПСИР). Показано, что метод позволяет удвоить порядок разнесения по сравнению с однозонными системами, обеспечивая выигрыш более 5 дБ по SNR при одинаковой аппаратурной сложности (числе АЦП). Подтверждена высокая устойчивость алгоритма к параметрическим неидеальностям конвертеров и его ортогональность к технологиям MIMO.

Результаты работы представляют интерес для проектирования энергоэффективных приёмных устройств нового поколения, где приоритет отдаётся надёжности и отказу от многоканальных гетеродинных трактов (при условии обеспечения входной аналоговой полосы тракта $[f_s, f_s+B]$) над максимальной спектральной эффективностью. Предлагаемая архитектура дополняет существующие стандарты (5G NR, Wi-Fi 6/7), предлагая альтернативный вектор развития для специализированных сегментов связи (IoT LPWAN, SatCom, Industrial Wireless), где аппаратные ресурсы и число антенн критичнее ширины занимаемого спектра. Дальнейшие исследования планируется направить на моделирование канального уровня с учётом помехоустойчивого кодирования и экспериментальную верификацию на платформах FPGA/RFSoC.

\vspace{0.3cm}
\noindent\textbf{Конфликт интересов и программное обеспечение.} Автор заявляет об отсутствии конфликта интересов. Численное моделирование и машинные тесты выполнены на языке Python с использованием открытых научных библиотек NumPy, SciPy и Matplotlib; исходный код полностью оригинален и доступен в открытом репозитории~\cite{github}.

\begin{thebibliography}{99}

\bibitem{abdelgader}Abdelgader A. M. S., Wu L. Design of a Multiorder OFDM Frequency Diversity Approach // International Journal of Antennas and Propagation. -- 2015. -- Vol. 2015. -- Art. ID 291625. -- 10 p. -- DOI: 10.1155/2015/291625.
\bibitem{adi_mt008}Analog Devices. MT-008 Tutorial: Converting Oscillator Phase Noise to Time Jitter [Электронный ресурс]. -- 2008. -- Режим доступа: \url{https://www.eetimes.com/understanding-the-effect-of-clock-jitter-on-high-speed-adcs-part-2-of-2/} (дата обращения: 20.11.2026).
\bibitem{black}Blackman R. B., Tukey J. W. The Measurement of Power Spectra from the Point of View of Communications Engineering. -- Dover Publications, 1958.
\bibitem{elbornsson}Elbornsson J., Gustafsson F., Eklund J. E. Blind adaptive equalization of mismatch errors in time-interleaved A/D converter systems // IEEE Transactions on Circuits and Systems I: Regular Papers. -- 2005. -- Vol. 52, No. 1. -- P. 151--159. -- DOI: 10.1109/TCSI.2004.840076.
\bibitem{fudge}Fudge G. L. et al. A Nyquist folding analog-to-information receiver // Proc. Asilomar Conf. Signals, Syst., Comput. -- 2008. -- P. 541--545. -- DOI: 10.1109/ACSSC.2008.5074477.
\bibitem{gadawe}Gadawe N. T., Hamad R. W., Qaddoori S. L. Efficient Implementation of Fractional Sampling Rate Conversion (SRC) on FPGA // Diagnostyka. -- 2025. -- Vol. 26, No. 3. -- P. 208854. -- DOI: 10.29354/diag/208854.
\bibitem{goldsmith}Goldsmith A. Wireless Communications. -- Cambridge University Press, 2005.
\bibitem{han}Han K., Kang S., Hong S. Sub-Nyquist Sampling OFDM Radar // IEEE Transactions on Radar Systems. -- 2023. -- Vol. 1. -- P. 669--680. -- DOI: 10.1109/TRS.2023.3333430.
\bibitem{harris}Harris F. J. Multirate Signal Processing for Communication Systems. -- Prentice Hall, 2004.
\bibitem{ieee_std}IEEE Std 1241-2023. Standard for Terminology and Test Methods for Analog-to-Digital Converters. -- IEEE, 2023.
\bibitem{murray}Murray M. J., Schermer R. T., Hart J., Murray J. B., Redding B. Deep learning enabled photonic Nyquist folding receiver for wideband RF spectral analysis // APL Photonics. -- 2025. -- Vol. 10, No. 2. -- P. 026112. -- DOI: 10.1063/5.0241958.
\bibitem{oppenheim}Oppenheim A. V., Schafer R. W. Discrete-Time Signal Processing. -- 3rd ed. -- Pearson, 2009.
\bibitem{proakis}Proakis J. G., Salehi M. Digital Communications. -- 5th ed. -- McGraw-Hill, 2008.
\bibitem{rajasekaran}Rajasekaran P. Bandpass Sampling: Lecture Notes for EE 4361 Introduction to Digital Signal Processing. -- Richardson: The University of Texas at Dallas, 2014. -- URL: \url{https://personal.utdallas.edu/~raja1/EE4361\%20Spring\%2014/Lecture\%20Notes/Bandpass\%20Sampling.pdf}.
\bibitem{rajeswaran}Rajeswaran A., et al. Time-Interleaved ADC Architectures and Calibration Techniques // Analog Integrated Circuits and Signal Processing. -- 2020.
\bibitem{razavi}Razavi B. RF Microelectronics. -- 2nd ed. -- Prentice Hall, 2012. -- ISBN 978-0-13-713473-1.
\bibitem{tropp}Tropp J. A., Gilbert A. C. Signal recovery from random measurements via orthogonal matching pursuit // IEEE Trans. Inf. Theory. -- 2007. -- Vol. 53, No. 12. -- P. 4655--4666.
\bibitem{tsai}Tsai S.-H., Wang H., Chen C.-L. Multi-channel equalization for time-interleaved analog-to-digital converters // IEEE Trans. Circuits Syst. II. -- 2007. -- Vol. 54, No. 10. -- P. 894--898.
\bibitem{vaidyanathan}Vaidyanathan P. P. Multirate Systems and Filter Banks. -- Prentice Hall, 1993.
\bibitem{vaughan}Vaughan R. G., Scott N. L., White D. R. The Theory of Bandpass Sampling // IEEE Transactions on Signal Processing. -- 1991. -- Vol. 39, No. 9. -- P. 1973--1984. -- DOI: 10.1109/78.134395.
\bibitem{walden}Walden R. H. Analog-to-digital converter survey and analysis // IEEE Journal on Selected Areas in Communications. -- 1999. -- Vol. 17, No. 4. -- P. 539--550.
\bibitem{github}Рычков Е. Н. Репозиторий с исходным кодом моделирования архитектуры ПСИР [Электронный ресурс]. -- 2026. -- Режим доступа: \url{https://github.com/Lennen/OFDM-Nyquist} (дата обращения: 20.11.2026).

\end{thebibliography}



\end{document}